\documentclass[12pt]{article}

% --------------------------------------------------
% Typography and encoding
% --------------------------------------------------
\usepackage[T1]{fontenc}
\usepackage[utf8]{inputenc}
\usepackage{newtxtext,newtxmath} % Times-style text and mathematics
\usepackage{microtype} % Improved spacing and justification

% --------------------------------------------------
% Page layout and spacing
% --------------------------------------------------
\usepackage[
margin=1in,
includefoot
]{geometry}

\usepackage{setspace}
\onehalfspacing

\setlength{\parindent}{1.5em}
\setlength{\parskip}{0pt}
\setlength{\emergencystretch}{3em}

% --------------------------------------------------
% Section formatting
% --------------------------------------------------
\usepackage{titlesec}

\titleformat{\section}
{\large\bfseries}
{\thesection.}
{0.5em}
{}

\titleformat{\subsection}
{\normalsize\bfseries}
{\thesubsection.}
{0.5em}
{}

\titleformat{\subsubsection}
{\normalsize\itshape}
{\thesubsubsection.}
{0.5em}
{}

\titlespacing*{\section}{0pt}{2.5ex plus 1ex minus 0.2ex}{1.2ex}
\titlespacing*{\subsection}{0pt}{2ex plus 0.8ex minus 0.2ex}{0.8ex}
\titlespacing*{\subsubsection}{0pt}{1.5ex plus 0.5ex minus 0.2ex}{0.5ex}

% --------------------------------------------------
% Figures and tables
% --------------------------------------------------
\usepackage{graphicx}
\usepackage{booktabs}
\usepackage{threeparttable}
\usepackage{caption}
\usepackage{subcaption}

\captionsetup{
font=small,
labelfont=bf,
labelsep=period,
justification=centering
}

% --------------------------------------------------
% Lists
% --------------------------------------------------
\usepackage{enumitem}
\setlist{
itemsep=0.25em,
topsep=0.5em,
parsep=0pt,
partopsep=0pt
}

\title{Model Sketch V1}
\author{José Zobaran}
\date{2026}

\begin{document}

\maketitle


\section{What do we want a model for?}
There are three linked questions we cannot answer from the data and econometrics only:
\begin{itemize}
  \item Decomposition of gains: redistribution versus aggregate increases
  \item From that, net program effect on production
  \item Accurate cost benfit analysis under no-program counterfactual
\end{itemize}
\vspace{1em}
A spatial equilibrium model could deliver these.

\section{Proposed Model}
Environment:
\begin{itemize}
  \item Investors with farm projects originate from a nonproductive region O
  \item They choose where to install their farm on any of N productive regions
  \item Moving costs are increasing in the distance from the region to O 
  \item Productivity is local and match specific
  \item Farms rent land and capital
  \item Land markets are locally determined; the capital market is nationally integrated  
  \item The program acts by making local capital more available
  \item Farms produce a single good at a fixed unit price
  \item Investors choses where to form a farm (or not to invest) and how much to produce
\end{itemize}
\vspace{1em}

So that the investor solves:
\begin{align}
\max_{K_i, L_{i,l}, l} \qquad A_l\eta_{i, l}K_i^\alpha L_{i, l}^\beta - r_l^k K_i - r_l^L L_{i, l} - m(d_l)
\end{align}
Given a 


Operationalization requires:
\begin{itemize}
  \item Moving costs; local productivity - timely proxies available
  \item Capital and land shares - 2007 agricultural census; IDB est. of 0.7 and 0.1 (Labor: 0.2)
  \item Capital supply elasticity - Harder; strong assumptions/parameter sweep may be needed
  \item Program impact on capital - Difficult to estimate due to data; could make assumptions based on program size
  \item Match specific productivity - extreme-value calibration
\end{itemize}
\vspace{1em}

Main possible issues and reasons:
\begin{itemize}
  \item Firm-focused; no welfare outcome - beyond the scope
  \item All land is productive; no role for extensive gains - weak evidence
  \item No labor input - industrial farmwork is not labor intensive
\end{itemize}
Moreover, I wanted to restrict and reduce the model as much as possible 
\end{document}